% ============================================================================
%  chapter_19_red_planet.tex
%  Chapter 19: The Red Planet
%  chem+ Applied Chemistry and Planetary Materials Science
%  VSEPR-SIM Theoretical Division
%  Compile: pdflatex (twice for refs / TOC)
% ============================================================================
\documentclass[11pt,a4paper]{article}

%% ── Encoding & fonts ──────────────────────────────────────────────────────────
\usepackage[utf8]{inputenc}
\usepackage[T1]{fontenc}
\usepackage{lmodern}
\usepackage{microtype}

%% ── Core packages ────────────────────────────────────────────────────────────
\usepackage{amsmath,amssymb,amsfonts}
\usepackage{mathtools}
\usepackage{bm}
\usepackage{geometry}
\usepackage{xcolor}
\usepackage{booktabs}
\usepackage{array}
\usepackage{tabularx}
\usepackage{enumitem}
\usepackage{fancyhdr}
\usepackage{titlesec}
\usepackage{hyperref}
\usepackage{float}
\usepackage{caption}
\usepackage{mdframed}
\usepackage{graphicx}

%% ── Page geometry ────────────────────────────────────────────────────────────
\geometry{left=2.8cm,right=2.8cm,top=3.2cm,bottom=3.0cm,headheight=15pt}

%% ── Hyperref ─────────────────────────────────────────────────────────────────
\hypersetup{
  colorlinks=true, linkcolor=black, citecolor=black, urlcolor=black,
  pdftitle={Chapter 19 -- The Red Planet},
  pdfsubject={chem+ Applied Chemistry: Mars Regolith and Iron Oxidation States}
}

%% ── Section numbering: 19.0 through 19.9 ────────────────────────────────────
\renewcommand{\thesection}{19.\arabic{section}}
\setcounter{section}{-1}
\renewcommand{\thesubsection}{\Alph{subsection}}

%% ── Section styling ──────────────────────────────────────────────────────────
\titleformat{\section}
  {\normalfont\large\bfseries}{\thesection.}{0.6em}{}
  [\vspace{-2pt}\rule{\textwidth}{0.4pt}]
\titleformat{\subsection}
  {\normalfont\normalsize\bfseries}{\thesubsection.}{0.5em}{}
\titleformat{\subsubsection}
  {\normalfont\normalsize\itshape}{\thesubsubsection.}{0.5em}{}

%% ── Header / footer ──────────────────────────────────────────────────────────
\pagestyle{fancy}
\fancyhf{}
\lhead{\footnotesize\textit{chem+}\;\textbar\;Chapter~19}
\rhead{\footnotesize\textit{The Red Planet}}
\lfoot{\footnotesize\today}
\rfoot{\footnotesize\thepage}
\renewcommand{\headrulewidth}{0.4pt}
\renewcommand{\footrulewidth}{0.3pt}

%% ── Chem+ identity bracket notation ─────────────────────────────────────────
%%
%%  \cid {species}{super}{sub}   full identity block:  [[species]]^super_sub
%%  \cidb{species}{sub}          subscript only:        [[species]]_sub
%%  \cidp{species}{super}        superscript only:      [[species]]^super
%%  \cidn{species}               bare brackets:         [[species]]
%%
\newcommand{\cid}[3]{\!\left[\!\left[\,#1\,\right]^{\!#2}_{\!#3}\!\right]}
\newcommand{\cidb}[2]{\!\left[\!\left[\,#1\,\right]_{\!#2}\!\right]}
\newcommand{\cidp}[2]{\!\left[\!\left[\,#1\,\right]^{\!#2}\!\right]}
\newcommand{\cidn}[1]{\!\left[\!\left[\,#1\,\right]\!\right]}

%% ── Named objects used throughout this chapter ───────────────────────────────
\newcommand{\Rmars}{\vec{R}_{\mathrm{Mars}}}   % regolith identity vector
\newcommand{\Fmag}{\vec{F}_{\mathrm{mag}}}     % magnetite state vector
\newcommand{\Rmat}{\mathcal{R}_{\mathrm{Mars}}}% regolith state matrix
\newcommand{\Proc}{\mathcal{P}}                % processing operator

%% ── Key-result box ───────────────────────────────────────────────────────────
\newmdenv[
  topline=true, bottomline=true, rightline=false, leftline=false,
  linewidth=1.2pt, linecolor=black,
  innertopmargin=8pt, innerbottommargin=8pt,
  innerleftmargin=4pt, innerrightmargin=4pt,
  skipabove=10pt, skipbelow=10pt
]{keybox}

%% ── Commentary / interjection box ────────────────────────────────────────────
\newmdenv[
  topline=true, bottomline=true, rightline=false, leftline=false,
  linewidth=1.2pt, linecolor=black,
  innertopmargin=8pt, innerbottommargin=8pt,
  innerleftmargin=0pt, innerrightmargin=0pt,
  skipabove=12pt, skipbelow=12pt
]{smlaw}

\captionsetup{font=small,labelfont=bf}

%% ═════════════════════════════════════════════════════════════════════════════
\begin{document}
%% ═════════════════════════════════════════════════════════════════════════════

%% ── Chapter opening ──────────────────────────────────────────────────────────
\thispagestyle{plain}
\vspace*{0.6cm}

\begin{center}
  \noindent\rule{\textwidth}{1.2pt}\vspace{0.8em}

  {\small\scshape Chapter~19}\vspace{0.3em}

  {\fontsize{22}{28}\selectfont\bfseries The Red Planet}

  \vspace{0.5em}
  \noindent\rule{\textwidth}{0.5pt}\vspace{0.5em}

  {\itshape
	Mars regolith as a chem+ system: iron oxidation states,
	contamination pathways,\\[0.2em]
	and the engineering chemistry of surface processing.
  }
\end{center}

\vspace{1.0em}
\noindent\rule{\textwidth}{0.4pt}
\vspace{0.8em}

\tableofcontents
\newpage

%%═══════════════════════════════════════════════════════════════════════════════
\section{Core Idea}
%%═══════════════════════════════════════════════════════════════════════════════

Mars is red because the surface is dominated by oxidised iron-bearing minerals.
In chem+ language, the visible planet is basically a planetary-scale
contamination layer:

\[
\cid{\mathrm{Mars~Regolith}}
	{\mathrm{Fe/O/Si/Al/Mg/Ca/Cl/S}}
	{s}
\]

with the red identity carried mainly by:

\begin{gather*}
\cid{\mathrm{Fe_2O_3}}{\mathrm{Fe{-}O{-}Fe}}{s}\\[6pt]
\cid{\mathrm{FeOOH}}{\mathrm{Fe{-}O{-}H}}{s}\\[6pt]
\cid{\mathrm{Fe_3O_4}}{\mathrm{Fe^{2+}/Fe^{3+}}}{s}
\end{gather*}

So the surface colour is not ``Mars is red,'' like a children's book written
by a committee.  It is:

\begin{keybox}
\[
\cidb{\mathrm{Fe}}{\text{reduced/rock},\;s}
\;\rightarrow\;
\cid{\mathrm{Fe^{3+}{-}O}}{\text{oxidised dust}}{s}
\]
\end{keybox}

The planet is red because iron has been converted into oxidised surface-state
solids.

%%═══════════════════════════════════════════════════════════════════════════════
\section{Red Planet Identity Vector}
%%═══════════════════════════════════════════════════════════════════════════════

Define the regolith identity vector:

\[
\Rmars
=
\begin{bmatrix}
\mathrm{SiO_2}\\[2pt]
\mathrm{Fe_2O_3}\\[2pt]
\mathrm{Al_2O_3}\\[2pt]
\mathrm{MgO}\\[2pt]
\mathrm{CaO}\\[2pt]
\mathrm{SO_4^{2-}}\\[2pt]
\mathrm{ClO_4^-}\\[2pt]
\mathrm{H_2O_{ads}}
\end{bmatrix}
\]

In compressed chem+ notation:

\[
\cid{\mathrm{Regolith}}
	{\mathrm{silicate/oxide/sulfate/perchlorate}}
	{s}
\]

Expanded:

\begin{align*}
\cidb{\mathrm{Regolith}}{\mathrm{Mars}}
={}&
  \cidb{\mathrm{Si{-}O{-}Si}}{s}
  +\cidb{\mathrm{Fe{-}O{-}Fe}}{s}
  +\cidb{\mathrm{Al{-}O{-}Si}}{s}\\
&+\cidb{\mathrm{Mg{-}O}}{s}
  +\cidb{\mathrm{Ca{-}SO_4}}{s}
  +\cidb{\mathrm{ClO_4^-}}{salt}
  +\cidb{\mathrm{H_2O}}{ads}
\end{align*}

The iron oxide piece is the red-state carrier:

\begin{keybox}
\[
\cid{\mathrm{Fe{-}O{-}Fe}}{\text{red chromophore}}{s}
\]
\end{keybox}

\begin{smlaw}
Yes, the planet's branding department is rust.  Space has no shame.
\end{smlaw}

%%═══════════════════════════════════════════════════════════════════════════════
\section{Red-State Formation Pathway}
%%═══════════════════════════════════════════════════════════════════════════════

A simplified oxidation path:

\[
\cidb{\mathrm{Fe}}{\text{metal/silicate},\;s}
\;+\;\cidb{\mathrm{O_2}}{atm}
\;+\;\cidb{\mathrm{H_2O}}{ads}
\;\rightarrow\;
\cidb{\mathrm{FeOOH}}{s}
\]

Then dehydration / ageing:

\[
2\,\cidb{\mathrm{FeOOH}}{s}
\;\rightarrow\;
\cidb{\mathrm{Fe_2O_3}}{s}
\;+\;\cidb{\mathrm{H_2O}}{g/ads}
\]

Chem+ flow:

\[
\cidp{\mathrm{Fe}}{\text{reduced}}_{\!\text{rock}}
\;\rightarrow\;
\cid{\mathrm{Fe{-}OH}}{\text{weathered}}{\text{surface}}
\;\rightarrow\;
\cid{\mathrm{Fe{-}O{-}Fe}}{\text{oxidised}}{\text{dust}}
\]

Main state transition:

\begin{keybox}
\[
\mathrm{Fe^0/Fe^{2+}}
\;\rightarrow\;
\mathrm{Fe^{3+}}
\]
\end{keybox}

Bond transition:

\begin{keybox}
\[
\cidb{\mathrm{Fe{-}Si/O}}{\text{rock}}
\;\rightarrow\;
\cidb{\mathrm{Fe{-}OH}}{\text{weathered}}
\;\rightarrow\;
\cidb{\mathrm{Fe{-}O{-}Fe}}{\text{oxide}}
\]
\end{keybox}

%%═══════════════════════════════════════════════════════════════════════════════
\section{The Red Planet Contamination Lesson}
%%═══════════════════════════════════════════════════════════════════════════════

The earlier acid-leach contamination pathway becomes important because Mars
regolith contains iron oxides everywhere.  The danger is not merely that iron
exists.  The danger is this conversion:

\begin{keybox}
\[
\cid{\mathrm{Fe}}{\text{immobile oxide}}{s}
\;\rightarrow\;
\cid{\mathrm{Fe}}{\text{mobile ferric}}{aq}
\]
\end{keybox}

On Mars, this matters for:

\begin{center}
\medskip
\begin{tabular}{@{}lll@{}}
$\cidn{\text{water extraction}}$   &
$\cidn{\text{oxygen production}}$  &
$\cidn{\text{regolith processing}}$\\[6pt]
$\cidn{\text{construction binders}}$ &
$\cidn{\text{metal extraction}}$   &
$\cidn{\text{sample purity}}$
\end{tabular}
\medskip
\end{center}

because any acidic or reactive processing can drag iron into the liquid phase
and ruin separation.

\begin{smlaw}
Humanity goes to Mars and immediately has to ask: ``How do we not dissolve
the entire red problem into our product stream?''
Charming species.  Very focused.
\end{smlaw}

%%═══════════════════════════════════════════════════════════════════════════════
\section{Revised Mars Contamination Pathway}
%%═══════════════════════════════════════════════════════════════════════════════

Using the earlier pathway, the full acid-dissolution sequence for a surface
iron oxide is:

\begin{align*}
&\cid{\mathrm{Fe_2O_3}}{\mathrm{Fe^{3+}{-}O^{2-}{-}Fe^{3+}}}{s}\\[8pt]
&\quad\xrightarrow{+\,\mathrm{H^+}}\;
  \cid{\mathrm{Fe{-}OH{-}Fe}}{\text{protonated}}{\text{surface}}\\[8pt]
&\quad\xrightarrow{+\,\mathrm{H^+}}\;
  \cid{\mathrm{Fe{-}OH_2^+}}{\text{leaving}}{\text{surface}}\\[8pt]
&\quad\xrightarrow{\text{cleavage}}\;
  \cid{\mathrm{Fe^{3+}}}{\text{mobile}}{aq}\\[8pt]
&\quad\xrightarrow{+\,\mathrm{H_2O}}\;
  \cid{\mathrm{Fe(H_2O)_6^{3+}}}{\text{aquo}}{aq}\\[8pt]
&\quad\xrightarrow{+\,\mathrm{SO_4^{2-}}}\;
  \cid{\mathrm{FeSO_4^+}}{\text{inner/outer sphere}}{aq}\\[8pt]
&\quad\xrightarrow{\text{cluster}}\;
  \cid{\mathrm{Fe{-}}\mu\mathrm{SO_4{-}Fe}}{\text{bridged}}{aq}
\end{align*}

Mars version, combining multiple dissolved iron species:

\begin{align*}
&\cid{\mathrm{Mars~dust}}{\mathrm{Fe_2O_3/FeOOH}}{s}
 \;+\;\cid{\mathrm{acidic~brine}}{\mathrm{H^+/SO_4^{2-}/Cl^-}}{aq}\\[6pt]
&\quad\rightarrow\;
  \cid{\mathrm{Fe^{3+}}}{\text{mobile}}{aq}
  \;+\;\cidb{\mathrm{FeSO_4^+}}{aq}
  \;+\;\cidb{\mathrm{FeCl^{2+}}}{aq}
  \;+\;\cidb{\mathrm{FeOH^{2+}}}{aq}
\end{align*}

So the Mars red-state problem becomes:

\begin{keybox}
\[
\text{red dust}
\;\rightarrow\;
\text{dissolved ferric contamination}
\;\rightarrow\;
\text{dirty brine chemistry}
\]
\end{keybox}

%%═══════════════════════════════════════════════════════════════════════════════
\section{Chem+ Bond-Pattern Library for Mars Iron}
%%═══════════════════════════════════════════════════════════════════════════════

%%─────────────────────────────────────────────────────────────────────────────
\subsection{Hematite-like Oxide}
%%─────────────────────────────────────────────────────────────────────────────

\[
\cid{\mathrm{Fe_2O_3}}{\mathrm{Fe^{3+}{-}O^{2-}{-}Fe^{3+}}}{s}
\]

Bond motif:
\begin{keybox}
\[\mathrm{Fe{-}O{-}Fe}\]
\end{keybox}

State:
\[
\cidb{\mathrm{Fe}}{+3,\;\text{oxide}}
\]

%%─────────────────────────────────────────────────────────────────────────────
\subsection{Hydroxide / Oxyhydroxide Surface}
%%─────────────────────────────────────────────────────────────────────────────

\[
\cid{\mathrm{FeOOH}}{\mathrm{Fe{-}O{-}H}}{s}
\]

Bond motif:
\begin{keybox}
\[\mathrm{Fe{-}OH}\]
\end{keybox}

Surface hydrated form:
\[
\cid{\mathrm{Fe{-}OH}}{\text{surface}}{\text{dust}}
\]

%%─────────────────────────────────────────────────────────────────────────────
\subsection{Magnetite-like Mixed Iron State}
%%─────────────────────────────────────────────────────────────────────────────

\[
\cid{\mathrm{Fe_3O_4}}{\mathrm{Fe^{2+}/Fe^{3+}}}{s}
\]

State vector:
\[
\Fmag
=
\begin{bmatrix}
\mathrm{Fe^{2+}}\\[2pt]
\mathrm{Fe^{3+}}\\[2pt]
\mathrm{O^{2-}}
\end{bmatrix}
\]

Mixed redox identity:
\[
\cid{\mathrm{Fe_3O_4}}{\text{mixed valence}}{s}
\]

This is important because mixed-valence iron can shift under chemical
processing:
\[
\mathrm{Fe^{2+}}
\;\rightleftharpoons\;
\mathrm{Fe^{3+}} + e^-
\]

%%─────────────────────────────────────────────────────────────────────────────
\subsection{Ferric Aquo Contaminant}
%%─────────────────────────────────────────────────────────────────────────────

\[
\cidb{\mathrm{Fe^{3+}}}{aq}
\;+\; 6\,\cidb{\mathrm{H_2O}}{l}
\;\rightarrow\;
\cidb{\mathrm{Fe(H_2O)_6^{3+}}}{aq}
\]

Chem+:
\[
\cidp{\mathrm{Fe}}{+3}_{\!aq}
\;\cdot\;
\cid{\mathrm{OH_2}}{\sigma\text{ donor}}{6}
\]

Bond motif:
\begin{keybox}
\[\mathrm{Fe{-}OH_2}\]
\end{keybox}

%%─────────────────────────────────────────────────────────────────────────────
\subsection{Ferric Sulfate Contaminant}
%%─────────────────────────────────────────────────────────────────────────────

\[
\cidb{\mathrm{Fe^{3+}}}{aq}
\;+\;\cidb{\mathrm{SO_4^{2-}}}{aq}
\;\rightleftharpoons\;
\cidb{\mathrm{FeSO_4^+}}{aq}
\]

Chem+:
\[
\cid{\mathrm{Fe{-}O{-}SO_3}}{\text{sulfate bound}}{aq}
\]

Bond motif:
\begin{keybox}
\[\mathrm{Fe{-}O{-}S}\]
\end{keybox}

%%─────────────────────────────────────────────────────────────────────────────
\subsection{Bridged Ferric Sulfate Cluster}
%%─────────────────────────────────────────────────────────────────────────────

\[
2\,\cidb{\mathrm{Fe^{3+}}}{aq}
\;+\;\cidb{\mathrm{SO_4^{2-}}}{aq}
\;\rightarrow\;
\cidb{\mathrm{Fe{-}O{-}SO_2{-}O{-}Fe}}{4+,\;aq}
\]

Chem+:
\[
\cid{\mathrm{Fe{-}}\mu\mathrm{SO_4{-}Fe}}{\text{bridged}}{aq}
\]

Bond motif:
\begin{keybox}
\[\mathrm{Fe{-}O{-}S{-}O{-}Fe}\]
\end{keybox}

\begin{smlaw}
This is where the dissolved iron stops being a clean ion and becomes a
coordination mess wearing a fake moustache.
\end{smlaw}

%%═══════════════════════════════════════════════════════════════════════════════
\section{Mars Processing Comparison}
%%═══════════════════════════════════════════════════════════════════════════════

%%─────────────────────────────────────────────────────────────────────────────
\subsection*{Acidic Route}
%%─────────────────────────────────────────────────────────────────────────────

\begin{align*}
&\cid{\mathrm{Mars~Regolith}}{\mathrm{Fe_2O_3/SiO_2/Al_2O_3}}{s}
 \;+\;\cidb{\mathrm{H^+}}{aq}\\[6pt]
&\quad\rightarrow\;
  \cidb{\mathrm{Fe^{3+}}}{aq}
  \;+\;\cidb{\mathrm{Al^{3+}}}{aq}
  \;+\;\cidb{\mathrm{Mg^{2+}}}{aq}
  \;+\;\cidb{\mathrm{Ca^{2+}}}{aq}
  \;+\;\cidb{\mathrm{SiO_2}}{\text{residue}}
\end{align*}

Problem:
\begin{keybox}
\[
\cidb{\mathrm{Fe_2O_3}}{s}
\;\rightarrow\;
\cidb{\mathrm{Fe^{3+}}}{aq}
\]
\end{keybox}

Acid makes extraction stronger but dirtier.

%%─────────────────────────────────────────────────────────────────────────────
\subsection*{Basic Route}
%%─────────────────────────────────────────────────────────────────────────────

\begin{align*}
\cidb{\mathrm{Mars~Regolith}}{s}
\;+\;\cidb{\mathrm{OH^-}}{aq}
\;\rightarrow\;
\cidb{\text{silicate/aluminate}}{aq}
\;+\;\cidb{\mathrm{Fe_2O_3}}{s/\text{residue}}
\end{align*}

Ideal selectivity:
\[
\cidb{\mathrm{Fe_2O_3}}{s}
\;\not\rightarrow\;
\cidb{\mathrm{Fe^{3+}}}{aq}
\]

\begin{keybox}
\[\text{basic leach keeps more iron locked as solid oxide}\]
\end{keybox}

\begin{smlaw}
The base route may be less aggressive, but it is cleaner for keeping red iron
out of solution.  Naturally, the slower path is the more civilised one.
Annoying but typical.
\end{smlaw}

%%═══════════════════════════════════════════════════════════════════════════════
\section{Regolith State Matrix}
%%═══════════════════════════════════════════════════════════════════════════════

Define the Mars regolith state matrix:

\[
\Rmat
=
\begin{bmatrix}
\mathrm{Fe_2O_3}   & \mathrm{FeOOH}       & \mathrm{Fe_3O_4}  \\[4pt]
\mathrm{SiO_2}     & \mathrm{Al_2O_3}     & \mathrm{MgO}      \\[4pt]
\mathrm{CaSO_4}    & \mathrm{NaCl/MgCl_2} & \mathrm{ClO_4^-}
\end{bmatrix}
\]

\noindent
Rows:
\[
\begin{bmatrix}
\text{redox / iron}\\[2pt]
\text{silicate structure}\\[2pt]
\text{salts / reactive contaminants}
\end{bmatrix}
\]

\noindent
Columns:
\[
\begin{bmatrix}
\text{oxide} & \text{hydrated / surface} & \text{reactive / mobile}
\end{bmatrix}
\]

A processing operation acts on this matrix:

\[
\Proc : \Rmat \;\rightarrow\; \Rmat'
\]

\medskip\noindent
Four canonical operators:

\begin{align*}
\Proc_{\text{acid}}    &:\;
  \cidb{\mathrm{Fe{-}O{-}Fe}}{s}
  \;\rightarrow\;
  \cidb{\mathrm{Fe^{3+}}}{aq}
\\[8pt]
\Proc_{\text{base}}    &:\;
  \cidb{\mathrm{Si{-}O{-}Al}}{s}
  \;\rightarrow\;
  \cidb{\text{silicate/aluminate}}{aq}
\\[8pt]
\Proc_{\text{thermal}} &:\;
  \cidb{\mathrm{H_2O}}{ads}
  \;\rightarrow\;
  \cidb{\mathrm{H_2O}}{g}
\\[8pt]
\Proc_{\text{redox}}   &:\;
  \cidb{\mathrm{Fe^{3+}}}{s}
  + 3e^-
  \;\rightarrow\;
  \cidb{\mathrm{Fe^0}}{s/l}
\end{align*}

%%═══════════════════════════════════════════════════════════════════════════════
\section{Red Planet Extraction Decision}
%%═══════════════════════════════════════════════════════════════════════════════

The main engineering question:

\begin{keybox}
\[\text{Do we want iron mobile, fixed, reduced, or separated?}\]
\end{keybox}

%%─────────────────────────────────────────────────────────────────────────────
\subsection*{Path 1 --- Preserve iron as solid contamination}
%%─────────────────────────────────────────────────────────────────────────────

\[
\cidb{\mathrm{Fe_2O_3}}{s}
\;\rightarrow\;
\cidb{\mathrm{Fe_2O_3}}{\text{residue}}
\]

Useful when extracting water, salts, or light elements.

%%─────────────────────────────────────────────────────────────────────────────
\subsection*{Path 2 --- Dissolve iron intentionally}
%%─────────────────────────────────────────────────────────────────────────────

\[
\cidb{\mathrm{Fe_2O_3}}{s}
\;+\; 6\,\cidb{\mathrm{H^+}}{aq}
\;\rightarrow\;
2\,\cidb{\mathrm{Fe^{3+}}}{aq}
\;+\; 3\,\cidb{\mathrm{H_2O}}{l}
\]

Useful for chemical analysis, but bad for clean product streams.

%%─────────────────────────────────────────────────────────────────────────────
\subsection*{Path 3 --- Reduce iron to metal}
%%─────────────────────────────────────────────────────────────────────────────

\[
\cidb{\mathrm{Fe_2O_3}}{s}
\;+\; 3\,\cidb{\mathrm{H_2}}{g}
\;\rightarrow\;
2\,\cidb{\mathrm{Fe}}{s}
\;+\; 3\,\cidb{\mathrm{H_2O}}{g}
\]

Chem+:
\[
\cidb{\mathrm{Fe^{3+}{-}O}}{s}
\;\rightarrow\;
\cidb{\mathrm{Fe^0}}{s}
\]

Useful for metallurgy and construction feedstock.

%%─────────────────────────────────────────────────────────────────────────────
\subsection*{Path 4 --- Convert oxide oxygen into useful oxygen}
%%─────────────────────────────────────────────────────────────────────────────

\[
\cidb{\mathrm{MO}}{s}
\;\rightarrow\;
\cidb{\mathrm{M}}{s/l}
\;+\;\cidb{\mathrm{O_2}}{g}
\]

General oxygen extraction:
\[
\cidb{\text{oxide}}{\text{regolith}}
\;\rightarrow\;
\cidn{\text{metal/silicate}}
\;+\;\cidn{\mathrm{O_2}}
\]

\begin{smlaw}
On Mars, oxygen extraction is not just life support.  It is propellant,
combustion, welding, processing, and not dying stupidly in a very expensive
desert.
\end{smlaw}

%%═══════════════════════════════════════════════════════════════════════════════
\section{The Red Planet as a Chem+ System}
%%═══════════════════════════════════════════════════════════════════════════════

Final chapter identity:

\begin{keybox}
\begin{align*}
\cid{\mathrm{Mars}}{\text{red}}{\text{planet}}
={}&\cidb{\text{silicate body}}{\text{rock}}\\
&+\;\cidb{\mathrm{Fe^{3+}{-}O}}{\text{dust}}\\
&+\;\cidb{\text{salts/perchlorates}}{\text{surface}}\\
&+\;\cidb{\mathrm{CO_2}}{\text{atm}}\\
&+\;\cidb{\mathrm{H_2O}}{\text{ice/ads}}
\end{align*}
\end{keybox}

The useful abstraction:

\[
\cidn{\text{Mars surface}}
\;=\;
\cidn{\text{structure}}
+\cidn{\text{oxidation}}
+\cidn{\text{salt}}
+\cidn{\text{dust}}
+\cidn{\text{volatile}}
\]

Or in one clean chem+ statement:

\begin{keybox}
\begin{align*}
\cidn{\text{Red Planet}}
={}&
  \cid{\mathrm{Fe{-}O{-}Fe}}{\text{colour/state}}{\text{dust}}\\
&+\;\cid{\mathrm{Si{-}O{-}Si}}{\text{structure}}{\text{rock}}\\
&+\;\cid{\mathrm{SO_4^{2-}/ClO_4^-}}{\text{reactive salts}}{\text{regolith}}\\
&+\;\cid{\mathrm{CO_2/H_2O}}{\text{volatile}}{\text{atm/ice}}
\end{align*}
\end{keybox}

\end{document}
